\chapter{Introduction to Category Theory}

Category theory abstracts the notion of internal structure and studies mathematical objects indirectly through their connections to other objects. In practical terms, instead of studying objects in isolation, category theory shifts the focus onto morphisms among objects, and instead of studying mathematical entities from a definitional stand point, it emphasizes the importance of how constructions relate to one another.

The rule-book of categorical thinking keeps us away from manipulating mathematical objects directly, it also offers a categorical lens as a tool to look at constructions in a holistic spirit. This lens abstracts away the details of mathematical constructions but enables a more general approach to  problems. This perspective unifies many areas of mathematics under a common language.

Category theory offers extreme generality and has great power when it comes to unifying seemingly unrelated fields. Which often misleads people into interpreting it as a unifying theory of mathematics that has no practical applications, when in reality it is more like a tool used to create blueprints that describe mathematical environments and their properties. It is also commonly used to give structure to patterns that repeat within different fields of mathematics. In some instances, category theory serves as a medium for transporting mathematical statements--such as theorems or lemmas--between domains. A categorical formulation of such results can often be translated to other contexts via functors or equivalences of categories.

One of the fields that benefits from the practical applications of categorical reasoning is computer science, deeply related to category theory through type theory, intuitionistic logic and natural deduction. We only mention CCC's and LCCC's as relevant constructions in the effort not to give away the purpose of the chapter. The unifying nature of categories brings us closer to the central focus of this work.

\newpage
\section{\centering Fundamentals of Category Theory}

Let us begin our dive into categories by defining what they are:

\begin{definition} A category $\mathcal{C}$ consists of:
  \begin{itemize}
  \item A collection of objects $\mathrm{Ob}(\mathcal{C})$
  \item A collection of arrows or morphisims where each arrow $f$ has a domain object $\mathrm{dom}(f)$ and a codomain object $\mathrm{cod}(f)$. These can be defined using $f : A \to B $ to signify $\mathrm{dom}(f) = A$ and $\mathrm{cod}(f) = B$
  \item A composition operator assigning to each pair of arrows $f$ and $g$, with $cod(f) = dom(g)$, a composite arrow $g \circ f : \mathrm{dom}(f) \to \mathrm{cod} (g)$ satisfying the following associative law:
    \[
      h \circ (g \circ f) = (h \circ g) \circ f
    \]
    where $f : A \to B, g : B \to C, h : C \to D$ are arrows of a category
  \item For each object $A$, an identity arrow $id_A : A \to A$ satisfying the following identity law:
    \[
      \mathrm{id_B} \circ f = f \quad f \circ \mathrm{id_A} = f 
    \]
    for any $f : A \to B$
  \end{itemize}
\end{definition}



\begin{remark}
Literature on the topic usually differentiates between small and large categories. Small categories are those whose collection of objects and whose collection of arrows are both sets. Large categories allow for collections of objects or arrows that are proper classes rather than sets. We will only be concerned with small categories.
\end{remark}

Various mathematical constructs fit the description of a category. For example, the category $\mathbf{Set}$ consists of sets as objects and functions as morphisms, where composition is ordinary function composition and identities are identity functions. Another example is the category $\mathbf{Grp}$, whose objects are groups and whose morphisms are group homomorphisms. Other examples are:


\begin{example} Let $\mathcal{C}$ denote a category, consisting of a set of objects, $\functionfont{Ob}(\mathcal{C})$, and a set of arrows, $\functionfont{Hom}(\mathcal{C})$, satisfying the category axioms.

\begin{itemize}
  \item
    \textbf{empty category}, denoted by $\mathbf{0}$, is the category with no objects and no morphisms.
    \[\text{Ob}(\mathbf{0}) = \emptyset, \quad \text{Hom}(\mathbf{0}) = \emptyset \]
  \item
    \textbf{terminal category} (or \textbf{point category}), denoted by $\mathbf{1}$, is the category with a single object and a single morphism, which must be the identity morphism on that object. Formally,
    \[ \text{Ob}(\mathbf{1}) = \{*\}, \quad \text{Hom}(\mathbf{1}) = \{\operatorname{id}_*\} \]
    where $\operatorname{id}_*: * \to *$ is the identity morphism on $*$.
  \item
    \textbf{category \(\mathbf{2}\)}, often representing the poset $A \le B$, is a category with two objects and three morphisms. Formally,
    \[ \text{Ob}(\mathbf{2}) = \{A, B\}, \quad \text{Hom}(\mathbf{2}) = \{\operatorname{id}_A, \operatorname{id}_B, f\} \]
    where $\operatorname{id}_A: A \to A$ and $\operatorname{id}_B: B \to B$ are identity morphisms, and $f: A \to B$ is a unique non-identity morphism. The compositions are defined as $f \circ \operatorname{id}_A = f$ and $\operatorname{id}_B \circ f = f$.
\end{itemize}

  
  The most elemental category is the category that has no elements and no arrows, it is commonly represented with a $0$. Category $1$ contains one object and one identity arrow, as required by the definition. Category $2$ contains three arrows, the identities and one arrow from $1$ to $2$. Category $3$ builds up on this as is observed in the following diagram:
  \[
    \begin{tikzcd}
      1 \arrow[loop right, "id_1"] & & & & & 2 \arrow[d] \arrow[loop right, "id_2"] \\
      1 \arrow[r] \arrow[loop left, "id_1"] & 2 \arrow[loop right, "id_2"] & & &1 \arrow[r] \arrow[ru] \arrow[loop left, "id_1"] & 3 \arrow[loop right, "id_3"]
    \end{tikzcd}
  \]
\end{example}
\begin{example} In this example we show that \textbf{Set} is a category where objects correspond to sets $A,B,C,\dots$ and arrows to functions.
  For $f: A \to B$ and $g: B \to C$, composition of functions is $(g \circ f)(a) = g(f(a))$, which is a function $ g \circ f : A \to C$.
  Composition is associative since for $h: C \to D$ we have that: \( h \circ (g \circ f)(a) = h(g(f(a))) = (h \circ g)(f(a)) = (h \circ g) \circ f(a) \).
  Lastly, each set $A$, the identity morphism is $\mathrm{id}_A(a) = a$, satisfying $f \circ \mathrm{id}_A = f = \mathrm{id}_B \circ f$.
  Therefore, $\mathbf{Set}$ is a category. \qed
\end{example}
\begin{example}
In this example we show that any partially ordered set $(P, \leq)$ can be viewed as a category $\mathbf{Pos}$ whose objects are elements of $P$ and there is a unique arrow $p \to q$ 
if and only if $p \leq q$.  
Composition is defined by transitivity: if $p \leq q$ and $q \leq r$, then $p \leq r$, 
so the composite $p \to r$ exists and is unique.  
Associativity holds automatically since there is at most one arrow between any two objects.  
Lastly, each $p \in P$ has an identity arrow $p \to p$ that corresponds to the reflexivity property of partial orders.
Therefore, every poset defines a category. \qed
\end{example}
%\orange{There exist several well known categories tied to mathematics and compite. Category of categories... Monoid}

\begin{example}
Let us consider a simple functional
programming language with the following types:
\[
  T ::= \mathrm{Nat} \mid \mathrm{Bool} \mid T \times T \mid T \to T
\]
The type system is closed under product and function type
constructors. A valid compound type is \( (\mathrm{Nat} \times \mathrm{Bool}) \to \mathrm{Nat} \).
A categorical interpretation of this scenario reveals types are objects and functions morphisms.

Identity \(\mathrm{id}_A : A \to A\) can be defined as \( \mathrm{id}_A = \lambda x\,{:}\,A . \, x  \), which is well-typed. Identity laws are verified trivially: \( \mathrm{id}_B \circ f = f = f \circ \mathrm{id}_A \).

Compositionality is defined as sequential application of functions that share domain and codomain through \(f : A \to B\) and \(g : B \to C\) like \(g \circ f : A \to C\) by \(g \circ f = \lambda x\,{:}\,A . \, g (f(x))\), which is well-typed by function application.

Composition is associative due to \(\beta\)-equivalence in lambda calculus:
\[
h \circ (g \circ f) = \lambda x{:}A . h (g (f(x))) = (h \circ g) \circ f
\]
Therefore, the language satisfies all axioms of a category.
\end{example}

\begin{example}
$\mathbf{Hask}$ is a category where the objects are Haskell types such as \small{\texttt{Int}, \texttt{Bool}, \texttt{String}}. The morphisms are total Haskell functions between those types such as  \small{\(\texttt{f::Int->Bool}\)}. For every type \small{\(\texttt{a}\)}, there is an identity morphism \small{\(\texttt{id::a->a}\)} that maps each value to itself.
Morphisms can be composed using \small{\(\texttt{(.)::(b->c)->(a->b)->a->c}\)},
meaning the output of one function becomes the input of another.
This composition is associative, and the identity morphisms act as neutral elements for composition.
These properties — types as objects, functions as morphisms, \small{\(\texttt{id}\)} as identity, and \texttt{\(\texttt{(.)}\)} as composition — make Hask a category.
\end{example}

\begin{table*}[h!]
\centering
\begin{tabular}{ll}
\hline
\textbf{Category} & \textbf{Description} \\
\hline
$\mathbf{Set}$      & Sets as objects, functions as morphisms. \\
$\mathbf{Grp}$      & Groups and group homomorphisms. \\
$\mathbf{Ring}$     & Rings and ring homomorphisms. \\
$\mathbf{Mod}_R$    & Modules over a ring $R$ and $R$-linear maps. \\
$\mathbf{Vect}_k$   & Vector spaces over field $k$ and linear maps. \\
$\mathbf{Top}$      & Topological spaces and continuous functions. \\
$\mathbf{Man}$      & Smooth manifolds with smooth maps. \\
$\mathbf{Hilb}$     & Hilbert spaces and bounded linear maps. \\
\hline
\end{tabular}
\caption{Key mathematical categories.}
\label{tab:math-categories}
\end{table*}

\begin{table*}[h!]
\centering
\begin{tabular}{ll}
\hline
\textbf{Category} & \textbf{Description} \\
  \hline
$\mathbf{Mon}$ & Monoids and monoid homomorphisms. \\
$\mathbf{CPO}$        & Complete partial orders and continuous maps. \\
$\mathbf{Dom}$        & Domains as posets with Scott-continuous maps. \\
$\mathbf{Rel}$        & Sets and binary relations as morphisms. \\
$\mathbf{Pfn}$        & Sets and partial functions as morphisms. \\
$\mathbf{Par}$        & Sets and partial maps (generalized morphisms). \\
$\mathbf{Kl}(T)$      & Kleisli category of a monad $T$ on a base category. \\
$\mathbf{Endo(Set)}$  & Endofunctors on the category of sets. \\
\hline
\end{tabular}
\caption{Key categories in computer science.}
\label{tab:cs-categories}
\end{table*}

%\begin{note}
%  It is ironic that the names of the better known categories are given by their objects and not their arrows.
%\end{note}

%\orange{to build some intuition on why categories relate to programs}
%\magenta{que alguien que controle corrija el siquiente ejemplo q es mas inventado que que los gamusinos }

Aside from defining categories from mathematical objects, categores can be built from other categories. One way to do this is via opposite categories.

\begin{definition} For a category $\mathcal{C}$, the opposite or dual category $\mathcal{C}^{op}$ is defined by:
\begin{itemize}
    \item $\mathrm{Ob}(\mathcal{C}^{op}) = \mathrm{Ob}(\mathcal{C})$,
    \item For $A,B \in \mathrm{Ob}(\mathcal{C})$, $\mathcal{C}^{op}(A,B) = \mathcal{C}(B,A)$,
    \item Composition is reversed: if $f : A \to B$ and $g : B \to C$ in $\mathcal{C}$, then
    \[
    f^{op} \circ g^{op} = (g \circ f)^{op}.
    \]
\end{itemize}
\end{definition}

\subsubsection{\centering Hom-set}

\begin{definition}
In a category $\mathcal{C}$, for any two objects $A, B \in \functionfont{Ob}(\mathcal{C})$, the hom-set from $A$ to $B$ is the set of all morphisms from $A$ to $B$. This set is denoted by $\functionfont{Hom}(A, B)$.
\[
    \functionfont{Hom}(A, B) = \{ f \mid f \colon A \to B \text{ is a morphism in } \mathcal{C} \}
\]
\end{definition}

\subsubsection{\centering Monos \& Epis}

\begin{definition}
  \label{def:monos-epis} Let \( g_1, g_2 : X \to A \) (resp. \( h_1, h_2 : B \to Y \)). An arrow \( f : A \to B \) in \( \mathcal{C} \) is a monomorphism (resp. epimorphism) if:
\[
  f \circ g_1 = f \circ g_2 \ \implies g_1 = g_2 \quad \text{(resp. } h_1 \circ f = h_2 \circ f \implies h_1 = h_2 \text{)}
\]

\end{definition}

\begin{note}
  Monomorphisims generalize the notion of injectivity, while epimorphisms are the categorical generalization of surjectivity.
\end{note}


\subsubsection{\centering Functors}
Categories provide a way of representing mathematical structures using only objects and arrows. Category theory takes itself very seriously, and since its focus is not on the objects themselves but on the relations between them, it is natural to consider functors, which are themselves morphisms in the category of categories. In this way, category theory can study itself.

With this perspective, the definition of a functor arises naturally as a pair of functions that map both objects and arrows from one category to another while preserving identities and composition. Formally:
\begin{definition} Let $\mathcal{C}$ and $\mathcal{D}$ be categories. A functor $F : \mathcal{C} \to \mathcal{D}$ consists of:
\begin{itemize}
\item A function that maps objects from $\mathcal{C}$ to $\mathcal{D}$:
  \[
    \begin{aligned}
      F : \mathrm{Ob}(\mathcal{C}) &\longrightarrow \mathrm{Ob}(\mathcal{D}) \\
      X &\longmapsto F(X)
    \end{aligned}
  \]
\item A function that maps arroes in $\mathcal{C}$ to $\mathcal{D}$:
  \[
    \begin{aligned}
      F : \mathrm{Hom}_{\mathcal{C}}(X,Y) &\longrightarrow \mathrm{Hom}_{\mathcal{D}}(F(X),F(Y)) \\
      f &\longmapsto F(f)
    \end{aligned}
  \]
\end{itemize}
such that the properties of a category are kept across the transformation:
\begin{itemize}
\item For every object $X \in \mathcal{C}$, $F(\mathrm{id}_X) = \mathrm{id}_{F(X)}$.
\item For all composable morphisms $f : X \to Y$ and $g : Y \to Z$ in $\mathcal{C}$,
  \[
    F(g \circ f) = F(g) \circ F(f).
  \]
\end{itemize}

The following diagram elegantly summarizes the latter definition:
\[
\begin{tikzcd}[row sep=2em, column sep=3em]
X \arrow[r, "f"] \arrow[d, mapsto, "F"'] & Y \arrow[d, mapsto, "F"] \\
F(X) \arrow[r, "F(f)"'] & F(Y)
\end{tikzcd}
\]
\end{definition}

\begin{example} Haskell captures the notion of functor using the following type class:
\begin{minted}{Haskell}
  class Functor f where
      fmap :: (a -> b) -> f a -> f b
      (<$) :: a -> f b -> f a
      (<$) = fmap . const
\end{minted}
There exist several instances for this class, among which we find:
\begin{minted}{Haskell}
  instance Functor [] where
      fmap = map
\end{minted}
That allows us to perform mapping operations on lists like:
\begin{minted}{shell-session}
  ghci> fmap (*2) [1,2,3]
  [2,4,6]
\end{minted}
\end{example}

\begin{example} Consider the categories of groups $\mathbf{Grp}$ and sets $\mathbf{Set}$. 
We define the functor \( U : \mathbf{Grp} \longrightarrow \mathbf{Set} \):
\begin{itemize}
    \item \textbf{On objects}: for each group $G \in \mathbf{Grp}$, set $U(G)$ to be the set of elements of $G$. 
    \item \textbf{On morphisms}: for each group homomorphism $f: G \to H$, we define $U(f) = f$, now considered as a function $U(G) \to U(H)$ between sets.
\end{itemize}
This mapping preserves identities and composition. Diagrammatically, we can represent the action of $U$ on a morphism $f : G \to H$ as:
\[
\begin{tikzcd}[row sep=2em, column sep=3em]
G \arrow[r, "f"] \arrow[d, mapsto, "U"'] & H \arrow[d, mapsto, "U"] \\
U(G) \arrow[r, "U(f)"'] & U(H)
\end{tikzcd}
\]
Since this functor ``forgets'' the structure its domain category we refer to it as forgetful functor.
\end{example}






\begin{example} A reformulation from example 2.1.2 in \cite{pierce1991basic} that defines a functor called List from the category $\mathbf{Set} \to \mathbf{Set}$ to illustrate functors from a programming perspective.
  
\begin{itemize}
\item For every set $S$, we define a new set:
\[
    \mathrm{List}(S) := \{ L \mid L \text{ is a finite list of belements of } S \}
\]
Where the order of elements matters, and duplicates are allowed.

\item For every function $f : S \to S'$, we define a new function
\begin{align*}
  \mathrm{List}(f) : \mathrm{List}(S) &\to \mathrm{List}(S') \\
  \mathrm{List}(f)([s_1, s_2, \dots, s_n]) &:= [f(s_1), f(s_2), \dots, f(s_n)].
\end{align*}
So, say for instance we have $S = \{ 1, 2 \}$ and $f(s) = 2 \cdot s $. Then, $S' = \{2, 4\}$, \( \mathrm{List}(\{1, 2\}) = \{ [], [1], [2], [1,1], [1,2], [2,1], [2,2], [1,1,1], \dots \} \), and \( \mathrm{List}(f)(\{2, 4\}) = \{ [], [2], [4], [2,2], [2,4], [4,2], [4,4], [2,2,2], \dots \} \).
\end{itemize}
To confirm that \(\mathrm{List}\) is a functor, it must satisfy the following two laws:
\begin{itemize}
\item \textbf{Identity Law:} For every set $S$ and every list $L \in \mathrm{List}(S)$: \(
\mathrm{List}(\mathrm{id}_S)(L) = L. \)
\item \textbf{Composition Law:} For any two functions $f : S \to S'$ and $g : S' \to S''$, and any list $L \in \mathrm{List}(S)$:
\(
\mathrm{List}(g \circ f)(L) = (\mathrm{List}(g) \circ \mathrm{List}(f))(L).
\)
\end{itemize}
\end{example}

\begin{note} \( (\mathrm{List}(S), ++, []) \) is a monoid, where $++$ is associative list concatenation and $[]$ is the identity, satisfying $ [] ++ L = L = L ++ []$.
\end{note}

\subsubsection{\centering Natural Transformations}
To define natural transformations we will apply the same reasoning scheme as we did
with functors. We had categories and we defined functors as morphisms between
categories, we do the same with natural transformations, that are nothing
but functor morphisms.
Suppose we have two functors $F, G : \mathcal{C} \to \mathcal{D}$. 
These represent two distinct ways of translating $\mathcal{C}$ into $\mathcal{D}$. 
A natural transformation between $F$ and $G$ then provides a systematic way 
to relate their outputs. Formally:

\begin{definition}

  Let $\mathcal{C}$ and $\mathcal{D}$ be categories, and let $F, G : \mathcal{C} \to \mathcal{D}$ be functors. A natural transformation \( \alpha : F \Rightarrow G \) between them, denoted:
\[
\begin{tikzcd}
\mathcal{C}
\arrow[r, bend left=50, "F"{name=U}]
\arrow[r, bend right=50, "G"'{name=D}]
& \mathcal{D}
\arrow[Rightarrow, from=U, to=D, "\alpha" description]
\end{tikzcd}
\]
is an assignment to every object $X \in \mathrm{Ob}(\mathcal{C})$ of a morphism \( \alpha_X : F(X) \to G(X) \text{ in } \mathcal{D} \) called the component of $\alpha$ at $X$, such that for every morphism $f: X \to Y$ in $\mathcal{C}$, the following diagram commutes:

\[
\begin{tikzcd}
F(X) \arrow[r, "\alpha_X"] \arrow[d, "F(f)"'] & G(X) \arrow[d, "G(f)"] \\
F(Y) \arrow[r, "\alpha_Y"] & G(Y)
\end{tikzcd}
\]
\end{definition}
\noindent A Haskell example that makes the diagram commute is:
\begin{minted}{shell-session}
  ghci> (reverse . (fmap toUpper)) "abc"
  "CBA"
  ghci> ((fmap toUpper) . reverse) "abc"
  "CBA"
\end{minted}
This example is one of the simplest of its kind since we are dealing with an endofunctor in $\mathbf{Hask}$.
\begin{definition} An endofunctor on a category $\mathcal{C}$ is a functor $ F : \mathcal{C} \to \mathcal{C} $
that maps objects and morphisms of $\mathcal{C}$ to objects and morphisms of the same category, preserving identities and composition.
\end{definition}

\begin{example} This example is in principle the same as the latter, but may help to grasp the idea of functor in Haskell:
\begin{minted}{Haskell}
  data BinTree a
      = Empty
      | Node a (BinTree a) (BinTree a)
      deriving (Show, Eq)

  instance Functor BinTree where
      fmap _ Empty = Empty
      fmap f (Node x l r) = Node (f x) (fmap f l) (fmap f r)

  reversePairs :: BinTree a -> BinTree a
  reversePairs Empty = Empty
  reversePairs (Node x l r) = Node x (reversePairs r) (reversePairs l)
  
  testNaturality :: (Eq b) => (a -> b) -> BinTree a -> Bool
  testNaturality f t =
      (reversePairs . fmap f) t == (fmap f . reversePairs) t
\end{minted}
We now can call the \texttt{testNaturality} function to assert whether the \texttt{f} we provide makes the diagram commute.
\begin{minted}{shell-session}
  ghci> f n = 2 * n
  ghci> testNaturality f t
  True
\end{minted}
In this case, any function that just operates on the keys will be a valid operation since it does not alter structure.
\end{example}
\begin{note}
  From this example we extract a design pattern. For a given data structure whether a graph, a tree or otherwise, an operation that modifies structure can be considered a natural transform whenever we only deal with functions that don't modify structure.
\end{note}
\begin{note}
In standard Haskell, we can only create endofunctors since the syntax only allows us to work with types.
\end{note}
Natural transformations, as well as functors can be composed as well, we study now the two ways there are of composing natural transformations:
\begin{remark}
Given functors $F, G, H : \mathcal{C} \to \mathcal{D}$, natural transformations $\alpha: F \Rightarrow G$ and $\beta: G \Rightarrow H$ can be composed:
\[
\beta \circ \alpha : F \Rightarrow H
\]
\[
\begin{tikzcd}
\mathcal{C}
\arrow[r, bend left=100, "F"{name=U}]
\arrow[r, ""'{name=V}]
\arrow[r, bend right=100, "H"'{name=D}]
& \mathcal{D}
\arrow[Rightarrow, from=U, to=V, "\alpha" description]
\arrow[Rightarrow, from=V, to=D, "\beta" description]
\end{tikzcd}
\]
with components \( (\beta \circ \alpha)_X = \beta_X \circ \alpha_X \) for each object $X \in \mathrm{Ob}(\mathcal{C})$.
\end{remark}

\begin{remark}
Given functors $F, G, H : \mathcal{C} \to \mathcal{D}$, natural transformations $\alpha: F \Rightarrow G$ and $\beta: G \Rightarrow H$ can be composed \emph{vertically}:
\[
\beta \circ \alpha : F \Rightarrow H
\]
\[
\begin{tikzcd}
\mathcal{C}
\arrow[r, bend left=100, "F"{name=U}]
\arrow[r, ""'{name=V}]
\arrow[r, bend right=100, "H"'{name=D}]
& \mathcal{D}
\arrow[Rightarrow, from=U, to=V, "\alpha" description]
\arrow[Rightarrow, from=V, to=D, "\beta" description]
\end{tikzcd}
\]
with components 
\[
(\beta \circ \alpha)_X = \beta_X \circ \alpha_X
\] 
for each object $X \in \mathrm{Ob}(\mathcal{C})$.

\medskip

\noindent
Given categories $\mathcal{C}, \mathcal{D}, \mathcal{E}$, functors $F, G : \mathcal{C} \to \mathcal{D}$ and $J, K : \mathcal{D} \to \mathcal{E}$, and natural transformations $\alpha: F \Rightarrow G$ and $\gamma: J \Rightarrow K$, we can define the \emph{horizontal composition}:
\[
\gamma \ast \alpha : J \circ F \Rightarrow K \circ G
\]
with components
\[
(\gamma \ast \alpha)_X = \gamma_{G(X)} \circ J(\alpha_X) = K(\alpha_X) \circ \gamma_{F(X)}
\]
for each $X \in \mathrm{Ob}(\mathcal{C})$.
\end{remark}
\begin{remark}
The collection of all functors $F : \mathcal{C} \to \mathcal{D}$ together with natural transformations between them forms the functor category $\mathbf{Fun}(\mathcal{C}, \mathcal{D}) \in \mathrm{Ob}(\mathbf{Cat})$.
\end{remark}
\subsubsection{\centering Universal Constructions}
\begin{definition}
Let $\mathcal{C}$ be a category and let $A,B \in \mathrm{Ob}(\mathcal{C})$.  
A product of $A$ and $B$ is an object $A \times B$ together with morphisms
\( \pi_1 : A \times B \to A, \pi_2 : A \times B \to B \) such that for every
object $X$ with morphisms $f : X \to A$ and $g : X \to B$, there exists a
unique morphism $\langle f,g \rangle : X \to A \times B$ making the diagram commute:
\[
\begin{tikzcd}
  & X \arrow[dl, "f"'] \arrow[dr, "g"] \arrow[d, dashrightarrow, "{\langle f,g \rangle}"] & \\
A & A \times B \arrow[l, "\pi_1"] \arrow[r, "\pi_2"'] & B
\end{tikzcd}
\]
\[
  \pi_1 \circ \langle f,g \rangle = f
  \quad
  \pi_2 \circ \langle f,g \rangle = g
\]
\end{definition}

\begin{note}
  In these diagrams, the dashed arrow indicates a unique morphism guaranteed by the universal property of the product or coproduct. That is, for any mediating maps as in the diagram, there exists exactly one morphism making the diagram commute.
\end{note}

\begin{definition}
Let $\mathcal{C}$ be a category and let $A,B \in \mathrm{Ob}(\mathcal{C})$.  
A coproduct of $A$ and $B$ is an object $A + B$ together with morphisms
\( \iota_1 : A \to A + B, \iota_2 : B \to A + B \) such that for every object
$X$ with morphisms $f : A \to X$ and $g : B \to X$, there exists a unique
morphism $[f,g] : A + B \to X$ making the diagram commute:
\[
\begin{tikzcd}
A \arrow[r, "\iota_1"] \arrow[dr, "f"'] & A + B \arrow[d, dashrightarrow, "{[f,g]}"] & B \arrow[l, "\iota_2"'] \arrow[dl, "g"] \\
& X &
\end{tikzcd}
\]
\[
[f,g] \circ \iota_1 = f
\quad
[f,g] \circ \iota_2 = g
\]
\end{definition}





